\section{Code development and benchmarking}
\label{sec:benchmarking}

In all previous NNPDF releases, including NNPDF3.0, PDF evolution and
deep-inelastic scattering were computed using the internal {\tt
  FKgenerator} code.  As discussed in Sect.~\ref{sec:expdata}, in
NNPDF3.1 PDF evolution is now performed using the public {\tt APFEL}
code. Deep-inelastic scattering and fixed-target Drell Yan are also
computed using {\tt APFEL}.
As far as perturbative evolution is concerned, {\tt APFEL} has been
extensively tested against other publicly available codes, such as
{\tt HOPPET}~\cite{Salam:2008qg} and {\tt QCDNUM}~\cite{Botje:2010ay}
for the PDF evolution and {\tt OpenQCDrad}~\cite{Alekhin:2013nda} for
the calculation of heavy-quark structure functions in the massive
scheme.

We have performed an extensive benchmarking of {\tt APFEL} and {\tt
  FKgenerator}, also involving deep-inelastic structure function (as
already mentioned in Ref.~\cite{Ball:2016neh}) and Drell-Yan cross-sections.
In the process of this benchmarking, two bugs were found in
the {\tt FKgenerator} implementation of the DIS structure functions
(one related to target-mass corrections, the other in the expressions
of the $\mathcal{O}(\alpha_s)$ charge-current massive coefficient
functions): we checked explicitly that none of them produced an effect
on NNPDF3.0 PDFs that could be distinguished from a statistical
fluctuation.

Representative results of the benchmarking of deep-inelastic structure
functions are shown in Fig.~\ref{fig:benchmarking}, where we show the
relative difference between {\tt FKgenerator} and {\tt APFEL}
implementation, using NNPDF3.0 as input PDF set, for the CHORUS
charged current neutrino-nucleus reduced cross-sections, the NMC
proton reduced cross-sections and neutral and charged current cross-sections
 from the H1 experiments from the HERA-II dataset.  In each
case, we show theoretical predictions calculations at LO and in the FONLL-A,
-B and -C general-mass schemes. The two codes are in good agreement, with
differences at most being at the 1\% level, typically much smaller.
%
This statement holds for all perturbative orders.

%%%%%%%%%%%%%%
\begin{figure}[t]
\centering
\includegraphics[width=0.46\textwidth]{images/plots/CHORUSNU.pdf}
\includegraphics[width=0.46\textwidth]{images/plots/NMC.pdf}
\includegraphics[width=0.46\textwidth]{images/plots/H1HERA2NCEM.pdf}
\includegraphics[width=0.46\textwidth]{images/plots/H1HERA2CCEP.pdf}
\caption{\small Relative difference in the DIS structure functions
  computed with {\tt FKgenerator} and {\tt APFEL}, using NNPDF3.0 as
  input PDF set, for the CHORUS charged current neutrino-nucleus
  reduced cross-sections, the NMC proton reduced cross-sections and
  neutral and charged current cross-sections from the H1 experiments
  from the HERA-II dataset. Datasets are as in Tab.~1 of
  Ref.~\cite{Ball:2014uwa}.  For each dataset, we compare the
  theoretical calculations at LO and in the FONLL-A, B and
  C~\cite{Forte:2010ta} heavy-quark mass schemes.
  \label{fig:benchmarking}}
\end{figure}
%%%%%%%%%%%%%%%%%%


The benchmarking of Drell-Yan cross-sections is illustrated in
Fig.~\ref{fig:benchmarking2}, where we show the relative difference
between the {\tt FKgenerator} and {\tt APFEL} calculations at LO, NLO
and NNLO for the E605 $pd$ and E866 $pp$ cross-sections, again using
NNPDF3.0 as input.  Again, differences are at most at the 2\% level
and typically much smaller. We have traced these residual differences
were to the fact that the coverage of the large-$x$ region was
sub-optimal in the {\tt FKgenerator} calculation, and is now improved
in {\tt APFEL} thanks to of a better choice of input $x$ grid.

%%%%%%%%%%%%%%
\begin{figure}[t]
\centering
\includegraphics[width=0.49\textwidth]{images/plots/60x_DYE605_fixed.pdf}
\includegraphics[width=0.49\textwidth]{images/plots/60x_DYE886R_P.pdf}
\caption{\small Same as Fig.~\ref{fig:benchmarking} for fixed-target
  Drell-Yan cross-sections: results are shown at LO, NLO and NNLO for
  the E605 $pA$ and E866 $pp$ cross-sections datasets, as given in
  Tab.~2 of Ref.~\cite{Ball:2014uwa}.
 \label{fig:benchmarking2}} 
\end{figure}
%%%%%%%%%%%%%%%%%%
